\section{Implications for infrastructure}\label{sec:implications}
% Section D of claims.md only. Every sentence states that it is a design implication, not a
% validated technology. P1-P4 = D4-D7 (D7 weak). D8 is in Related Work.
% No \claim inside the table (marginpar in a float): the claims of each row are in the
% sentence that introduces the table. \todo{verify} in "current mechanism" wherever the
% source is not verified on the primary (docs/VERIFICHE-chiuse.md).
% Not written (candidate, NOT in claims.md): fixed per-request price (pay-per-crawl)
% against a cost that changes sign. Enters only with a D row and a primary source.
Everything in this section is a design implication that we draw from the measurements; none
of it is a validated technology.\claim{D3}

\paragraph{Classification during the transition.}
Our data imply that classification remains necessary: at the knee of the testbed,
blocking the exhaustive class by its identity lies on the frontier of served work and
human latency (\cref{sec:r-frontier}).\claim{D2}
They also imply that classification is not necessarily sufficient.\claim{D3}
A policy that knows only that traffic is agentic does not necessarily hold enough
information to estimate its marginal impact, since under the same label the marginal
cost ranged from $-0.035$ to $+0.438$ origin requests per request
(\cref{sec:r-sign}).\claim{D1}

\paragraph{Three quantities beside the label.}
As design implications, three quantities would complement the class label.\claim{D4}\claim{D5}\claim{D6}
The first implication, and the one with the strongest support, is that infrastructure
should be able to observe or derive the width of the working set of a class, since that
width alone changed the sign of the marginal cost.\claim{D4}\claim{A1}
The implication asks for a quantity that request logs already contain: we computed it
over a short window for the clients of our honeypot (\cref{sec:m-honeypot}).\claim{D4}\claim{C3}
The second implication is that it should know the marginal cost of the class in the
current state of the shared cache, because at the operating point of the frontier one
blocked exhaustive request frees about 0.82 origin requests and one agentic request about
0.10.\claim{D5}\claim{A5}
The third implication is that it should know the effect of a class on the other
workloads, because the human p99 rose by 37.0\,ms while the human miss ratio changed only
by a factor of 1.06 (\cref{sec:m-human}).\claim{D6}\claim{B5}
A fourth implication, weaker because our evidence for it is only indirect, concerns the
state of the shared cache and the pressure on resources.\claim{D7}

\paragraph{From evidence to properties.}
\Cref{tab:bridge} restates these implications as four properties, P1 to P4, and sets
each one against the evidence that motivates it, the mechanism in use today and where
that mechanism falls short; it is a statement of implications, not an
architecture.\claim{D4}\claim{D5}\claim{D6}\claim{D7}

\begin{table}[t]
\caption{Design implications, not validated technology: what infrastructure would need to
know (P1--P4 = claims D4--D7; P4 is weakly supported), the evidence behind each, the
mechanism in use today and where it falls short for the evidence we have.}
\label{tab:bridge}
\centering\scriptsize
\renewcommand{\arraystretch}{1.15}
\begin{tabularx}{\textwidth}{@{}l>{\raggedright\arraybackslash}X
  >{\raggedright\arraybackslash}X>{\raggedright\arraybackslash}X
  >{\raggedright\arraybackslash}X@{}}
\toprule
 & Evidence & What infrastructure would need to know & Mechanism in use today &
   Where it falls short \\
\midrule
P1 & Same class and volumes, marginal from $-0.035$ to $+0.438$ as only the working set
  widens (\cref{sec:r-sign}) &
  Width of the class's working set over a window &
  Identity: declared or verified bot categories with per-category
  actions~\cite{cloudflare2026options}; signed requests~\cite{ietf-webbotauth}\todo{verify} &
  The label is the same at both ends of the range \\
\addlinespace
P2 & Blocking frees ${\sim}0.82$ (exhaustive) and ${\sim}0.10$ (agentic) origin
  requests per request (\cref{sec:m-block}) &
  Marginal cost of the class in the current cache state &
  Cache partitioning by the gradient of hit or miss ratio in allocated
  memory~\cite{qureshi2006ucp,cidon2016cliffhanger} &
  Varies allocation, not the composition of the traffic that shares the cache \\
\addlinespace
P3 & Human p99 $+37.0$\,ms with human miss ratio $\times 1.06$ (\cref{sec:m-human}) &
  Effect of a class on the other workloads &
  Hit ratio of human traffic as the measure of harm~\cite{zhang2025,cloudflare2026cache};
  a separate cache layer for AI traffic~\cite{cloudflare2026cache} &
  The harm arrives as origin load, which the human hit ratio barely shows \\
\addlinespace
P4 (weak) & Indirect only &
  Shared state and pressure on resources &
  Load- or latency-driven admission control (Breakwater, DAGOR, Protego, TopFull)
  \todo{verify} &
  Not measured in this paper \\
\bottomrule
\end{tabularx}
\end{table}

% Phase 2 of the simulator (docs/PREREG-simulatore-fase2.md) tests whether P1-P3 hold
% beyond LRU and the synthetic sessions. Write only after its RISULTATO document and an S
% row in claims.md; SIMULATO, never MISURATO.
\todo{simulator phase 2: do P1--P3 hold beyond this cache policy and generator?}

\paragraph{An open question.}
To our knowledge, the question these implications open has not been posed in this form:
whether admission, placement and priority can move from deciding on identity to deciding
on state and cost.\claim{D1}\claim{D3}
% Section E rows: "current architecture inadequate" and "new protocol or architecture
% needed", both NON DIMOSTRATO.
We do not claim that current infrastructure is inadequate or that a new protocol or
architecture is needed; our data do not show either.
